% El hilo de la materia: clase por clase
\section{El hilo de la materia: clase por clase}\label{cap:02}\label{sec:hilo}

\textbf{Material docente.} El profesor publicó en el campus un resumen escrito de cada una de
sus diez clases de la cursada 2025 (segundo cuatrimestre 2025: la clase 1 es del 21 de agosto
de 2025 y el material se llama ``Clase 2025.2.2''); el campus los titula ``1er Cuatri 2026'',
pero acá se citan como ``cursada 2025''. Leídos de corrido muestran que la materia no
es una lista de temas sueltos sino un recorrido con tres ideas que vuelven una y otra vez.
Cuando una clase parece desconectada de la anterior, casi siempre es porque cambió la
\emph{cara} desde la que se mira la misma idea.

\begin{idea}[Las tres ideas que se repiten]
\begin{enumerate}
  \item \textbf{Transformaciones, operaciones y coordenadas} son tres caras de lo mismo. Una
  transformación mueve puntos; una operación (componer, multiplicar matrices, hacer un
  producto interno) es la cuenta que la representa; una coordenada es lo que queda cuando se
  pregunta ``hasta dónde llegué con esta transformación''. En palabras del profesor:
  ``una de las cosas importantes en esta materia es entender la relación entre
  transformaciones, operaciones y coordenadas'' (resumen de clase 2).
  \item \textbf{Ir a lo chico para volver a lo grande.} Una rotación grande es muchas
  rotaciones chiquitas; la derivada de una transformación da vectores tangentes a las
  trayectorias; la exponencial rearma el movimiento completo. La misma jugada reaparece con
  el gradiente, con las potencias de una matriz y con la primera forma fundamental de una
  superficie.
  \item \textbf{La composición es una operación entre transformaciones}, y una
  transformación se identifica con la trayectoria que genera al aplicarla repetidas veces.
  Por eso importa tanto que la composición no sea conmutativa: sin conmutatividad no se
  pueden armar líneas coordenadas.
\end{enumerate}
\end{idea}

Antes de la clase 1 va la Unidad 0 del sitio de la materia, ``fundamentos y lenguaje de
la cátedra'' (subsección~\ref{sec:desaprender}): no es una clase del profesor sino
el encuadre del curso con el que conviene entrar a la tabla. El sitio tiene tres capas: el
hilo conductor (las 20 unidades interactivas en marimo + Python, construidas ladrillo a
ladrillo en el orden de esta tabla), el mapa de conceptos, para repasar, y este apunte
formal en \LaTeX{} con las demostraciones (apéndice~\ref{ap:notebooks}).

La tabla ubica cada clase con la idea eje dicha con las palabras del resumen del profesor,
el hecho geométrico que está detrás de la fórmula, y la sección de este apunte donde se
desarrolla. Regla del apunte: una fórmula sin su hecho geométrico es magia.

{\small
\begin{longtable}{>{\raggedright\arraybackslash}p{20mm}>{\raggedright\arraybackslash}p{60mm}>{\raggedright\arraybackslash}p{44mm}>{\raggedright\arraybackslash}p{28mm}}
\toprule
\textbf{Clase} & \textbf{Idea eje (palabras del profesor)} & \textbf{Hecho geométrico} &
\textbf{En este apunte}\\
\midrule
\endfirsthead
\toprule
\textbf{Clase} & \textbf{Idea eje (palabras del profesor)} & \textbf{Hecho geométrico} &
\textbf{En este apunte}\\
\midrule
\endhead
\bottomrule
\endlastfoot
1. Rotaciones &
``Una rotación implica dejar quieto un punto, mover otro, y en base a esos acomodar todos
los demás''; se puede porque ``había una estructura de algo que se conservaba (distancias y
ángulos)''. Complejos, matrices ortogonales, rotaciones 3D, no conmutatividad, ángulos
chicos, exponencial. &
Rotación = conservar distancias alrededor de un punto fijo; cada punto recorre una
circunferencia y multiplicar por \(i\) es girar un cuarto de vuelta. &
\S\ref{sec:rotar}, \S\ref{sec:rot-plano}, \S\ref{sec:rot-complejos}, \S\ref{sec:rot-3d},
\S\ref{sec:eje-exp}, ap.~\ref{ap:rot3d}, ap.~\ref{ap:infinitesimales}\\
2. Más rotaciones, composición y líneas coordenadas &
``En general es necesario para construir un sistema de coordenadas que haya conmutatividad
a la hora de componer movimientos por coordenadas.'' El problema de los vasos; rotaciones
sobre la cáscara de la esfera hacia latitud y longitud. &
Componer en distinto orden lleva a lugares distintos; las trayectorias de un movimiento
repetido son las líneas coordenadas (meridianos y paralelos). &
\S\ref{sec:orden-importa}, \S\ref{sec:sistemas-marcos}, \S\ref{sec:curvilineas}\\
3. Transformaciones afines &
``Una afinidad conserva paralelas las paralelas''; con reflexiones se construyen
rotaciones; ``un tema insistente de la clase fue la composición como una operación entre
transformaciones''. Dos escrituras: en \(\R^2\) y en coordenadas homogéneas. &
Afinidad = los paralelogramos siguen siendo paralelogramos; dos reflexiones en rectas
secantes son una rotación; rotar y trasladar repetidamente dibuja círculos alrededor del
punto fijo. &
\S\ref{sec:afinidades-primera}, \S\ref{sec:reflejar}, \S\ref{sec:componer-centros},
\S\ref{sec:homogeneas-primera}, \S\ref{sec:bloque-euclideas}--\S\ref{sec:afinidades-construidas}\\
4. Espacios lineales, transformaciones lineales y formas lineales &
``Las formas lineales, funciones que van de un espacio vectorial a un escalar''; ``se
comportan igual que las ecuaciones de las rectas o los planos''; ``una forma lineal se
transforma con una transformación lineal del lado contrario que un vector''. &
Forma lineal = una regla graduada tirada en el plano: sus curvas de nivel son rectas
paralelas equiespaciadas; la matriz es una pila de reglas. &
\S\ref{sec:transf-plano}, \S\ref{sec:formas-lineales}\\
5. Subespacios y producto cuña &
``El producto cuña, una operación relacionada al determinante que sirve para armar formas
lineales''; ``construir una transformación lineal que tuviera un vector especial (o varios)
en el núcleo''. &
Producto cuña = área o volumen orientado; vale cero cuando los vectores están alineados;
una transformación con núcleo aplasta una dirección entera a un punto. &
\S\ref{sec:vectorial}, \S\ref{sec:determinante}, \S\ref{sec:exterior}, \S\ref{sec:nucleo}\\
6. Subespacios invariantes y autovectores &
``\((A-\lambda\,Id)\cdot\vec v=0\): hay un vector que está en el núcleo de
\(A-\lambda\,Id\)''; ``los subespacios invariantes de la transformación lineal de
\(\R^3\) correspondían a los puntos fijos en \(\mathbb{P}\R^2\)''. &
Autovector = dirección que cae sobre sí misma; en homogéneas es un punto fijo del plano;
el cero siempre queda quieto. &
\S\ref{sec:autovalores}, \S\ref{sec:autovalores-geom}, \S\ref{sec:homografia-fijos}\\
7. Transformaciones lineales II: subespacios y SVD &
``Encontrar una `base privilegiada' de la transformación''; ``siempre existe una base
ortogonal del espacio de entrada que se conserva ortogonal a la salida''. Polya y las
estrategias para resolver problemas. &
Diagonalizar = mirar la transformación desde sus direcciones fijas (estiramientos
axiales); SVD = la circunferencia unidad se vuelve elipse, con ejes ortogonales antes y
después. &
\S\ref{sec:autovalores-geom}, \S\ref{sec:svd} (anticipado), ap.~\ref{ap:notebooks}\\
8. Cambios de base, autovectores y Marimo &
``El cambio de base se construye con una matriz que tiene en sus columnas los vectores de
una base expresada en la otra''; ``no parece haber un conjunto de vectores privilegiado
a priori''; potencias de una matriz y trayectorias. &
Cambio de base = la misma flecha medida con otra grilla; potencia real de una matriz =
trayectoria continua de un punto bajo la transformación. &
\S\ref{sec:cambio-base}, ap.~\ref{ap:notebooks}\\
9. Repaso de análisis (con álgebra de contrabando) &
``La derivada es la mejor aproximación lineal a una función en determinado punto'';
``dato clave: el gradiente es una forma lineal, no un vector''. Hessiana como forma
cuadrática. &
Derivada = el plano tangente; gradiente = forma lineal cuya recta de nivel cero es la
tangente, perpendicular a la curva de nivel. &
\S\ref{sec:plano-superficie}, \S\ref{sec:variedades}, \S\ref{sec:metrica},
\S\ref{sec:formas-lineales}\\
10. Inversas &
``La inversa es un elemento de un conjunto que, al operar con otro, resulta en el elemento
neutro respecto de esa operación''; núcleo, imagen, co-imagen y espacio inaccesible;
pseudoinversa de Moore--Penrose. &
Inversa = deshacer el movimiento; una transformación que pierde dimensiones aplasta y no se
deshace del todo; la pseudoinversa elige el punto de la imagen más cercano. &
\S\ref{sec:pseudoinversa} (anticipado), \S\ref{cap:cuatro-problemas}, \S\ref{sec:sistemas-ecuaciones}\\
\end{longtable}
}

\begin{atencion}[Segundo cuatrimestre 2026: otro orden, mismo hilo]
El curso del segundo cuatrimestre 2026 sigue otro orden: rotaciones \(\to\) guía de
reflexiones y similaridades \(\to\) afinidades y coordenadas homogéneas \(\to\)
transformaciones proyectivas \(\to\) productos, determinante y autovectores. Las tres ideas
son las mismas y el apunte sirve para ambos recorridos: la tabla anterior dice dónde está
cada tema, y las cajas ``En palabras del profesor'' de las secciones ``La idea'' de cada capítulo marcan a qué
clase de la cursada 2025 corresponde cada sección.
\end{atencion}

